% ===========================================================================
%  DSCSP: Diffusion-Stego with Conditional Sampling Paths
%  for Robust and Undetectable Image Steganography — IEEE TIFS Journal Format
% ===========================================================================
\documentclass[journal]{IEEEtran}

% ─── Packages ──────────────────────────────────────────────────────────────
\usepackage{amsmath, amssymb, amsfonts}
\usepackage{graphicx}
\usepackage{booktabs}
\usepackage{hyperref}
\usepackage{algorithm}
\usepackage{algorithmic}
\usepackage{multirow}
\usepackage{enumitem}
\usepackage{caption}
\usepackage{subcaption}
\usepackage{xcolor}
\usepackage{bm}
\usepackage{microtype}
\usepackage{cite}

\hypersetup{
    colorlinks=true,
    linkcolor=blue!70!black,
    citecolor=green!50!black,
    urlcolor=blue!60!black,
}

% ─── Math shortcuts ────────────────────────────────────────────────────────
\newcommand{\R}{\mathbb{R}}
\newcommand{\E}{\mathbb{E}}
\newcommand{\Lcal}{\mathcal{L}}
\newcommand{\Acal}{\mathcal{A}}
\newcommand{\Ncal}{\mathcal{N}}
\newcommand{\Ucal}{\mathcal{U}}
\DeclareMathOperator{\tamark}{tanh}
\DeclareMathOperator{\sigmoid}{sigmoid}
\DeclareMathOperator{\BCE}{BCE}

% ===========================================================================
\title{DSCSP: Diffusion-Stego with Conditional Sampling Paths for Robust and Undetectable Image Steganography}

\author{
    \IEEEauthorblockN{Authors Omitted for Review}
    \IEEEauthorblockA{}
}

\begin{document}
\maketitle

% ===========================================================================
%  ABSTRACT
% ===========================================================================
\begin{abstract}
We present DSCSP (Diffusion-Stego with Conditional Sampling Paths), a steganographic framework that hides secret messages within images produced by pre-trained diffusion models.
A compact learned \emph{Controller} network translates payload bits into spatial perturbation patterns applied to the output of a frozen DDIM sampling chain, while a spatially-aligned \emph{Extractor} recovers the payload from the resulting stego image.
Unlike prior diffusion-based methods that modify the noise vector or fine-tune the model, DSCSP keeps the generative backbone completely frozen, enabling efficient training on consumer hardware.

While the base system achieves near-perfect bit recovery and high visual fidelity, initial stego images are trivially detectable by learned steganalysis classifiers (AUC\,=\,1.0).
To address this, we introduce an \emph{adversarial training} regime with a Spatial Rich Model (SRM)--enhanced discriminator that forces the Controller to produce statistically indistinguishable perturbations.
We further propose a \emph{five-phase curriculum} that progressively introduces single attacks, adversarial pressure, and \emph{combined multi-attack} compositions, along with an optional \emph{redundant encoding} scheme for crop robustness.

Experiments on CelebA-HQ at $256\times256$ resolution demonstrate that DSCSP~v2 embeds 256 bits with BER\,$<$\,0.001, PSNR\,=\,40.9\,dB, and SSIM\,=\,0.962, while substantially improving robustness against noise (BER\,$<$\,0.002), blur (BER\,$<$\,0.002), brightness/contrast (BER\,$<$\,0.006), and their sequential compositions.
Adversarial training drives the discriminator to equilibrium (accuracy\,=\,0.50) and reduces external steganalysis AUC from 1.0 to 0.999, though further architectural changes are needed to achieve full statistical undetectability.
A comprehensive ablation study across seven architectural variants validates the contribution of each design component.
\end{abstract}

\begin{IEEEkeywords}
    Steganography, diffusion models, adversarial training, deep learning, image security, robustness, steganalysis
\end{IEEEkeywords}

% ===========================================================================
%  I. INTRODUCTION
% ===========================================================================
\section{Introduction}
\label{sec:intro}

For decades, steganography has centered on tweaking pixel values or transform coefficients in existing cover images so that the changes stay below the threshold of human perception~\cite{cox2007digital}.
The rise of deep generative models has shifted this picture: in \emph{generative steganography}, the cover image is created from scratch with the secret already baked in, which sidesteps the need for a pre-existing cover and naturally provides plausible deniability.

Diffusion models~\cite{ho2020denoising, song2020denoising} now produce images rivaling real photographs in quality, which makes them appealing as the generative backbone for steganographic systems.
Current diffusion-based steganography work can be grouped into two camps.
\emph{Noise-projection methods}, such as Diffusion-Stego~\cite{wei2023diffusionstego}, encode the message into the starting noise vector $\mathbf{z}_T$ and then run the standard reverse process---an elegant training-free idea, but one that provides limited defense against post-processing.
\emph{Model-adaptation methods}~\cite{peng2023stegodiffusion, zhang2024ldstega} modify or extend the diffusion model weights, which is expensive and can degrade the model's original generation quality.

A critical shortcoming shared by most high-capacity steganographic methods is \emph{statistical detectability}: trained classifiers can often distinguish stego images from clean ones with near-perfect accuracy~\cite{boroumand2019deep}, undermining the fundamental premise of steganography.

We propose \textbf{DSCSP} (Diffusion-Stego with Conditional Sampling Paths), which treats the \emph{sampling trajectory} of a frozen pre-trained diffusion model as an information channel.
A compact \emph{Controller} network maps payload bits to spatial perturbation patterns applied to the frozen DDIM~\cite{song2020denoising} output, and a lightweight \emph{Extractor} recovers the hidden bits.
To address detectability, we introduce an \emph{adversarial training} framework with a Spatial Rich Model (SRM)~\cite{fridrich2012rich}--enhanced discriminator, forcing the Controller to produce perturbations that are statistically indistinguishable from clean image characteristics.
A \emph{five-phase curriculum} progressively introduces single attacks, adversarial pressure, and combined multi-attack compositions.

Our contributions are as follows:
\begin{itemize}[itemsep=2pt, topsep=2pt]
    \item We introduce \textbf{sampling-path encoding}, a novel steganographic channel for diffusion models that maps bits to learned spatial perturbation patterns applied to frozen DDIM outputs, achieving 256-bit payloads with near-zero BER, high PSNR, and strong robustness.

    \item We propose an \textbf{SRM-enhanced adversarial discriminator} with fixed high-pass filter preprocessing and one-sided label smoothing that drives the internal discriminator to equilibrium (accuracy\,=\,0.50) and reduces external steganalysis AUC from 1.0 to 0.999.

    \item We design a \textbf{five-phase training curriculum} that progressively introduces clean convergence, single attacks, adversarial pressure, full training, and combined multi-attack compositions for comprehensive robustness.

    \item We develop a \textbf{combined attack framework} that composes sequential multi-attack pipelines (e.g., JPEG $\to$ noise, crop $\to$ JPEG) during both training and evaluation, reflecting realistic transmission scenarios.

    \item We present a \textbf{comprehensive ablation study} across seven architectural variants that isolates the contribution of orthogonality regularization, curriculum scheduling, LPIPS loss, attack training, spatial alignment, and adversarial training.

    \item We provide \textbf{redundant encoding} with majority-vote decoding for improved crop robustness, reducing crop BER from $\sim$0.47 to $\sim$0.29.
\end{itemize}


% ===========================================================================
%  II. RELATED WORK
% ===========================================================================
\section{Related Work}
\label{sec:related}

\subsection{Classical Image Steganography}

Traditional steganography methods embed messages by modifying pixel values~\cite{cox2007digital} or transform-domain coefficients of existing cover images.
LSB replacement~\cite{chan2004hiding} modifies the least significant bits of pixel values but is easily detectable by statistical steganalysis.
More sophisticated approaches such as HUGO~\cite{pevny2010using}, WOW~\cite{holub2012designing}, and S-UNIWARD~\cite{holub2014universal} minimize a distortion function to reduce detectability, but remain fundamentally limited by the cover image distribution.

\subsection{Deep Learning-Based Steganography}

Neural network-based approaches learn the encoding and decoding functions jointly.
HiDDeN~\cite{zhu2018hidden} pairs an encoder and decoder with a discriminator and a differentiable JPEG proxy.
SteganoGAN~\cite{zhang2019steganogan} takes a GAN-based route, folding the message directly into the generation process.
These learned approaches push capacity well beyond classical methods, though their output quality is limited by the generative backbone.

\subsection{Diffusion-Based Steganography}

Diffusion-Stego~\cite{wei2023diffusionstego} projects the secret message into the initial noise vector through an orthogonal projection scheme---training-free but vulnerable to spatial transformations.
StegoDiffusion~\cite{peng2023stegodiffusion} exploits invertible diffusion architectures.
LDStega~\cite{zhang2024ldstega} works within the latent space of latent diffusion models.
DERO~\cite{zhao2024dero} trains a message encoder--decoder jointly with the diffusion model for robustness.
Our work departs from all of these: DSCSP attaches a small learned perturbation controller to the \emph{output} of a fully frozen diffusion model, keeping the generative backbone untouched and the training footprint small.

\subsection{Adversarial Training for Undetectability}

GAN-based steganography~\cite{zhang2019steganogan} and adversarial example literature~\cite{goodfellow2014generative} have demonstrated that adversarial training can force generators to produce outputs that fool discriminators.
In steganalysis, Spatial Rich Models (SRM)~\cite{fridrich2012rich} use hand-crafted high-pass filters to expose steganographic artifacts.
We integrate SRM preprocessing into a learned discriminator and train adversarially against it, combining the pattern-detection strengths of classical steganalysis with the adaptability of deep learning.


% ===========================================================================
%  III. METHODOLOGY
% ===========================================================================
\section{Methodology}
\label{sec:method}

\subsection{System Overview}
\label{sec:overview}

DSCSP consists of three core components: (1)~a frozen pre-trained diffusion model $D_\theta$ that generates high-quality base images, (2)~a trainable Controller $C_\phi$ that maps payload bits to a spatial perturbation signal, and (3)~a trainable Extractor $E_\psi$ that recovers the payload from the stego image.
In the v2 system, we add (4)~a trainable SRM-enhanced Discriminator $\mathcal{D}_\omega$ for adversarial undetectability training.

The encoding process works as follows.
Given a payload $\mathbf{b} \in \{0,1\}^{N_b}$ with $N_b = 256$ bits:
\begin{enumerate}[itemsep=1pt, topsep=2pt]
    \item The Controller computes a perturbation: $\bm{\delta} = C_\phi(\mathbf{b}) \in \R^{3 \times H \times W}$.
    \item The diffusion model generates a clean image via DDIM sampling: $\mathbf{x}_\text{clean} = D_\theta(\mathbf{z}_T)$ under \texttt{no\_grad}.
    \item The stego image is formed: $\mathbf{x}_s = \text{clamp}(\mathbf{x}_\text{clean} + \bm{\delta},\, -1,\, 1)$.
\end{enumerate}
For extraction, the Extractor produces logits: $\hat{\mathbf{b}} = \sigmoid(E_\psi(\mathbf{x}_s))$.
Bits are decoded by thresholding at $0.5$.

The gradient path is:
\begin{equation}
    \nabla_\phi \Lcal = \frac{\partial \Lcal}{\partial \hat{\mathbf{b}}} \cdot \frac{\partial \hat{\mathbf{b}}}{\partial \mathbf{x}_s} \cdot \frac{\partial \mathbf{x}_s}{\partial \bm{\delta}} \cdot \frac{\partial \bm{\delta}}{\partial \phi}
    \label{eq:grad_path}
\end{equation}
so gradients never propagate through the heavy diffusion UNet.

\subsection{Controller Network}
\label{sec:controller}

The Controller takes a 256-bit payload $\mathbf{b} \in \{0,1\}^{256}$ and produces a spatial perturbation map $\bm{\delta} \in \R^{3 \times 256 \times 256}$ through multi-bit spatial tiling.

\subsubsection{Spatial Grid}
We divide the $256 \times 256$ image into an $8 \times 8$ grid where each cell covers $32 \times 32$ pixels, producing $N_c = 64$ cells.
With $K = 4$ bits per cell, the grid accommodates $N_c \times K = 256$ bits.

\subsubsection{Learned Patterns}
Each cell $j \in \{1, \ldots, 64\}$ has $K = 4$ independent learned pattern tensors:
\begin{equation}
    \mathbf{P}_j = \{\mathbf{p}_{j,k}\}_{k=1}^{K}, \quad \mathbf{p}_{j,k} \in \R^{3 \times 32 \times 32}
    \label{eq:patterns}
\end{equation}

\subsubsection{Forward Pass}
Given bits $\mathbf{b} = (b_1, \ldots, b_{256})$:
\begin{align}
    &\text{Group:} \quad \mathbf{b}^{(j)} = (b_{4j-3}, \ldots, b_{4j}) \in \{0,1\}^4 \label{eq:group} \\
    &\text{Signs:} \quad s_{j,k} = 2 b^{(j)}_k - 1 \in \{-1, +1\} \label{eq:signs} \\
    &\text{Cell pattern:} \quad \mathbf{c}_j = \sum_{k=1}^{K} s_{j,k} \cdot \mathbf{p}_{j,k} \label{eq:cell} \\
    &\text{Tile:} \quad \bm{\delta}_\text{raw} = \text{tile}(\{\mathbf{c}_j\}_{j=1}^{64}) \in \R^{3 \times 256 \times 256} \label{eq:tile} \\
    &\text{Constrain:} \quad \bm{\delta} = \tanh(\bm{\delta}_\text{raw}) \cdot \alpha_\text{max} \label{eq:constrain}
\end{align}
where $\alpha_\text{max} = 0.3$ is the maximum perturbation magnitude.
The Controller has approximately \textbf{786K parameters} ($64 \times 4 \times 3 \times 32 \times 32$).

\subsection{Extractor Network}
\label{sec:extractor}

We employ a ResNet-18~\cite{he2016deep} backbone (without pre-training) with the final pooling and FC layers removed.
For a $256 \times 256$ input, the backbone's stride of 32 produces feature maps of shape $(B, 512, 8, 8)$, naturally aligning with the Controller's $8\times8$ grid.
A $1\times1$ convolutional head reduces 512 channels to $K = 4$:
\begin{equation}
    \mathbf{f} = \text{Conv}_{1\times1}^{512 \to 4}(\text{ResNet18}(\mathbf{x}_s)) \in \R^{B \times 4 \times 8 \times 8}
    \label{eq:extractor_head}
\end{equation}
The feature map is permuted from $(B, K, 8, 8)$ to $(B, 8, 8, K)$ before flattening to align with the Controller's bit layout.
The Extractor has approximately \textbf{11.2M trainable parameters}.

\subsection{Frozen Diffusion Backbone}
\label{sec:diffusion}

We use \texttt{google/ddpm-celebahq-256}, a pre-trained DDPM~\cite{ho2020denoising}, generating $256\times256$ face images via 20-step deterministic DDIM~\cite{song2020denoising} sampling.
The UNet (114M parameters) is loaded in FP16 and fully frozen.
Rather than running the diffusion model at every training step, we pre-generate a pool of $M = 128$ clean images and cache them in GPU memory, refreshing every 2{,}000 steps.
After pool generation, the UNet is offloaded to CPU to free GPU memory.

\subsection{Loss Functions}
\label{sec:losses}

The total training objective is:
\begin{equation}
    \Lcal = \Lcal_\text{msg} + \lambda_p \Lcal_\text{perc} + \lambda_c \Lcal_\text{ctrl} + \lambda_o \Lcal_\text{ortho} + \rho(t) \cdot \lambda_a \Lcal_\text{adv}
    \label{eq:total_loss}
\end{equation}
where $\rho(t)$ is the adversarial ramp function (Section~\ref{sec:adversarial}).

\subsubsection{Message Reconstruction Loss}
Binary cross-entropy with logits:
\begin{equation}
    \Lcal_\text{msg} = \BCE(\hat{\mathbf{b}}_\text{logits},\, \mathbf{b})
    \label{eq:loss_msg}
\end{equation}

\subsubsection{Perceptual Fidelity Loss}
LPIPS~\cite{zhang2018unreasonable} distance using a frozen VGG backbone:
\begin{equation}
    \Lcal_\text{perc} = \text{LPIPS}_\text{VGG}(\mathbf{x}_s,\, \mathbf{x}_\text{clean})
    \label{eq:loss_perc}
\end{equation}

\subsubsection{Controller Regularization}
Mean squared perturbation magnitude:
\begin{equation}
    \Lcal_\text{ctrl} = \frac{1}{3HW} \|\bm{\delta}\|_F^2
    \label{eq:loss_ctrl}
\end{equation}

\subsubsection{Orthogonality Regularization}
To prevent pattern collapse, we penalize deviation of the Gram matrix from identity:
\begin{equation}
    \Lcal_\text{ortho} = \frac{1}{N_c} \sum_{j=1}^{N_c} \left\| \mathbf{G}_j - \mathbf{I}_K \right\|_F^2
    \label{eq:loss_ortho}
\end{equation}
where $\mathbf{G}_j \in \R^{K \times K}$ is the Gram matrix of flattened patterns in cell $j$.

\subsubsection{Adversarial Loss}
The Controller's adversarial loss encourages the discriminator to classify stego images as clean:
\begin{equation}
    \Lcal_\text{adv} = -\log(\mathcal{D}_\omega(\mathbf{x}_s))
    \label{eq:loss_adv}
\end{equation}
This term is weighted by the adversarial ramp $\rho(t)$ which increases from 0 to 1 over the course of Phase~3 (Section~\ref{sec:curriculum}).

Loss weights: $\lambda_p = 0.1$, $\lambda_c = 0.0$, $\lambda_o = 0.001$, $\lambda_a = 0.1$.

\subsection{Adversarial Discriminator}
\label{sec:adversarial}

To reduce statistical detectability, we introduce an adversarial discriminator $\mathcal{D}_\omega$ that is trained to distinguish clean from stego images, while the Controller is trained to fool it.

\subsubsection{SRM Preprocessing}
Following classical steganalysis~\cite{fridrich2012rich}, the discriminator begins with a fixed (non-trainable) Spatial Rich Model (SRM) preprocessing layer.
Three high-pass filters---1st, 2nd, and 3rd order finite differences---are applied independently to each RGB channel, producing $3 \times 3 = 9$ residual feature maps:
\begin{equation}
    \mathbf{r}_i = \mathbf{K}_i \ast \mathbf{x}, \quad i \in \{1, \ldots, 9\}
    \label{eq:srm}
\end{equation}
where $\mathbf{K}_1, \mathbf{K}_2, \mathbf{K}_3$ are the SRM kernels:
\begin{equation}
    \mathbf{K}_1 = \begin{bmatrix} 0 & 0 & 0 \\ 0 & -1 & 1 \\ 0 & 0 & 0 \end{bmatrix}, \;
    \mathbf{K}_2 = \begin{bmatrix} 0 & 0 & 0 \\ 1 & -2 & 1 \\ 0 & 0 & 0 \end{bmatrix}, \;
    \mathbf{K}_3 = \begin{bmatrix} 0 & 0 & 0 \\ -1 & 3 & -3 \\ 1 & 0 & 0 \end{bmatrix}
    \label{eq:srm_kernels}
\end{equation}
These fixed filters suppress image content and amplify steganographic noise, making it easier for subsequent layers to detect perturbation patterns.

\subsubsection{Discriminator Architecture}
After SRM preprocessing, a 5-block strided CNN processes the 9-channel residual map:
\begin{equation}
    \mathcal{D}_\omega: \R^{9 \times H \times W} \to [0, 1]
\end{equation}
Each block applies Conv2d ($4\times4$, stride 2) $\to$ BatchNorm $\to$ LeakyReLU(0.2), with channel progression $9 \to 32 \to 64 \to 128 \to 256 \to 512$.
The output is globally average-pooled and mapped to a probability via a linear layer with sigmoid activation.
The discriminator has approximately \textbf{2.76M parameters}.

\subsubsection{Training Strategy}
We adopt alternating optimization: at each training step, the discriminator is updated first on a batch of clean and stego image pairs, then the generator (Controller + Extractor) is updated with the adversarial loss included.
The discriminator uses binary cross-entropy with \emph{one-sided label smoothing}---clean targets are set to 0.9 instead of 1.0, while stego targets remain at 0.1---which prevents the discriminator from becoming overconfident and stabilizes adversarial dynamics:
\begin{equation}
    \Lcal_\text{disc} = -[0.9 \cdot \log \mathcal{D}_\omega(\mathbf{x}_\text{clean}) + \log(1 - \mathcal{D}_\omega(\mathbf{x}_s) + 0.1)]
    \label{eq:disc_loss}
\end{equation}

\subsection{Combined Attack Framework}
\label{sec:combined}

Real-world image processing often involves multiple sequential transformations (e.g.,\ a social media platform might resize, compress, and add a watermark).
To handle this, we introduce a \emph{combined attack} framework that composes two or more differentiable attacks sequentially:
\begin{equation}
    \mathbf{x}' = \Acal_2(\Acal_1(\mathbf{x}_s))
\end{equation}

During training, with probability $p_\text{combined} = 0.3$ (of attacked batches), the attack layer selects $n = 2$ attacks uniformly at random from the five available types and applies them in sequence.
During evaluation, we test five fixed combined pipelines: JPEG$\to$Noise, JPEG$\to$Blur, Crop$\to$JPEG, Noise$\to$BrightnessContrast, and JPEG$\to$Noise$\to$Blur.

\subsection{Redundant Encoding}
\label{sec:redundant}

To improve crop robustness, we provide an optional redundant encoding scheme.
A \emph{RedundantCoder} maps $M$ effective message bits to $N = M \times r$ physical encoding slots (default $M = 128$, $r = 2$, $N = 256$), ensuring each message bit is encoded in $r$ non-adjacent spatial cells:
\begin{equation}
    \text{Encode}: \{0,1\}^M \to \{0,1\}^N, \quad
    \text{Decode}: [0,1]^N \to \{0,1\}^M
\end{equation}
Decoding applies \emph{majority voting} over the $r$ replicas of each bit, providing resilience when a subset of spatial cells is corrupted by cropping.

\subsection{Five-Phase Training Curriculum}
\label{sec:curriculum}

Training follows a five-phase curriculum with 30{,}000 total steps:

\begin{itemize}[itemsep=2pt, topsep=2pt]
    \item \textbf{Phase~1} (steps 0--3{,}000): \emph{Clean convergence.} No attacks, no adversarial training. The Controller and Extractor learn the basic encode--decode mapping.

    \item \textbf{Phase~2} (steps 3{,}000--5{,}000): \emph{Attack ramp-in.} Single-attack probability ramps linearly from 0 to the curriculum level. No adversarial training yet.

    \item \textbf{Phase~3} (steps 5{,}000--10{,}000): \emph{Adversarial ramp-in.} The discriminator begins training and the adversarial loss weight $\rho(t)$ ramps from 0 to 1. Attack probability continues ramping.

    \item \textbf{Phase~4} (steps 10{,}000--15{,}000): \emph{Full training.} Maximum attack probability (0.7), full adversarial weight. All components at full strength.

    \item \textbf{Phase~5} (steps 15{,}000--30{,}000): \emph{Combined attacks.} 30\% of attacked batches use sequential multi-attack compositions. The system learns to handle realistic composite distortions.
\end{itemize}

The attack probability schedule is:
\begin{equation}
    p_\text{atk}(t) = \begin{cases}
        0 & \text{if } t < 3{,}000 \\
        p_\text{max} \cdot \frac{t - 3{,}000}{7{,}000} & \text{if } 3{,}000 \leq t < 10{,}000 \\
        p_\text{max} & \text{if } t \geq 10{,}000
    \end{cases}
    \label{eq:curriculum}
\end{equation}
with $p_\text{max} = 0.7$.

The adversarial ramp is:
\begin{equation}
    \rho(t) = \begin{cases}
        0 & \text{if } t < 5{,}000 \\
        \frac{t - 5{,}000}{5{,}000} & \text{if } 5{,}000 \leq t < 10{,}000 \\
        1 & \text{if } t \geq 10{,}000
    \end{cases}
    \label{eq:adv_ramp}
\end{equation}

\subsection{Differentiable Attack Layer}
\label{sec:attacks}

We implement five differentiable attacks that preserve gradient flow:

\subsubsection{Gaussian Noise}
Additive zero-mean Gaussian noise with $\sigma \sim \Ucal(0.01, 0.05)$:
\begin{equation}
    \mathbf{x}' = \text{clamp}(\mathbf{x} + \bm{\epsilon},\, -1, 1), \quad \bm{\epsilon} \sim \Ncal(\mathbf{0}, \sigma^2 \mathbf{I})
\end{equation}

\subsubsection{Gaussian Blur}
Depthwise convolution with a $5\times5$ Gaussian kernel ($\sigma \sim \Ucal(0.5, 2.0)$).

\subsubsection{Brightness/Contrast}
Affine intensity transformation with $c \sim \Ucal(0.8, 1.2)$, $\beta \sim \Ucal(-0.1, 0.1)$:
\begin{equation}
    \mathbf{x}' = \text{clamp}(c \cdot \mathbf{x} + \beta,\, -1, 1)
\end{equation}

\subsubsection{Crop and Resize}
Random spatial crop with ratio $r \sim \Ucal(0.7, 0.95)$ followed by bilinear interpolation.

\subsubsection{Differentiable JPEG (STE)}
Full JPEG simulation with Straight-Through Estimator~\cite{bengio2013estimating} for quantization:
\begin{equation}
    \hat{d} = d + \left(\text{round}\!\left(\frac{d}{q}\right) \cdot q - d\right)_\text{detach}
    \label{eq:ste}
\end{equation}

\subsubsection{Attack Wrapper}
During training, the \texttt{AttackLayer} randomly selects one attack (or a combined sequence) per batch with probability $p_\text{atk}$, or passes the image through unchanged.

\subsection{Optimization}
\label{sec:optimization}

Optimization uses Adam~\cite{kingma2014adam} ($\beta_1 = 0.9$, $\beta_2 = 0.999$) with the Controller (no weight decay) and Extractor ($10^{-4}$ weight decay) in separate parameter groups, both starting at $5 \times 10^{-4}$ with cosine annealing to $10^{-5}$.
The discriminator uses a separate Adam optimizer with learning rate $2 \times 10^{-4}$ and weight decay $10^{-4}$.
Gradients are clipped at max norm 5.0.

\begin{algorithm}[t]
\caption{DSCSP Encoding}
\label{alg:encoding}
\begin{algorithmic}[1]
\REQUIRE Payload $\mathbf{b} \in \{0,1\}^{256}$, frozen $D_\theta$, trained $C_\phi$
\ENSURE Stego image $\mathbf{x}_s$
\STATE $\bm{\delta} \leftarrow C_\phi(\mathbf{b})$
\STATE $\mathbf{x}_\text{clean} \leftarrow D_\theta(\mathbf{z}_T)$ with \texttt{no\_grad}
\STATE $\mathbf{x}_s \leftarrow \text{clamp}(\mathbf{x}_\text{clean} + \bm{\delta},\, -1,\, 1)$
\RETURN $\mathbf{x}_s$
\end{algorithmic}
\end{algorithm}

\begin{algorithm}[t]
\caption{DSCSP v2 Training Step}
\label{alg:training}
\begin{algorithmic}[1]
\REQUIRE Batch of clean images $\mathbf{X}$, random bits $\mathbf{B}$
\STATE \textit{// Discriminator update}
\STATE $\bm{\delta} \leftarrow C_\phi(\mathbf{B})$; \; $\mathbf{X}_s \leftarrow \text{clamp}(\mathbf{X} + \bm{\delta}, -1, 1)$
\STATE $\Lcal_\text{disc} \leftarrow \BCE(\mathcal{D}_\omega(\mathbf{X}), 0.9) + \BCE(\mathcal{D}_\omega(\mathbf{X}_s.\text{detach}()), 0.1)$
\STATE Update $\omega$ via $\nabla_\omega \Lcal_\text{disc}$
\STATE \textit{// Generator update}
\STATE $\mathbf{X}'_s \leftarrow \text{AttackLayer}(\mathbf{X}_s, p_\text{atk}(t))$
\STATE $\hat{\mathbf{B}} \leftarrow E_\psi(\mathbf{X}'_s)$
\STATE $\Lcal \leftarrow \Lcal_\text{msg} + \lambda_p \Lcal_\text{perc} + \lambda_o \Lcal_\text{ortho} + \rho(t) \lambda_a \Lcal_\text{adv}$
\STATE Update $\phi, \psi$ via $\nabla_{\phi,\psi} \Lcal$
\end{algorithmic}
\end{algorithm}


% ===========================================================================
%  IV. EXPERIMENTS
% ===========================================================================
\section{Experiments}
\label{sec:experiments}

\subsection{Experimental Setup}
\label{sec:setup}

\subsubsection{Hardware}
All experiments are conducted on an NVIDIA RTX 3060 Laptop GPU with 6\,GB VRAM running PyTorch 2.5.1 with CUDA.

\subsubsection{Diffusion Model}
We use the pre-trained \texttt{google/ddpm-celebahq-256} DDPM from HuggingFace, generating $256\times256$ face images via 20-step DDIM sampling.
The UNet has approximately 114M parameters, all frozen in FP16.

\subsubsection{Configuration}
Table~\ref{tab:hyperparams} summarizes all hyperparameters for v2 training.

\begin{table}[t]
\centering
\caption{Training hyperparameters for DSCSP v2.}
\label{tab:hyperparams}
\small
\begin{tabular}{@{}llr@{}}
\toprule
\textbf{Category} & \textbf{Parameter} & \textbf{Value} \\
\midrule
\multirow{4}{*}{Model}
    & Payload capacity & 256 bits \\
    & Bits per cell & 4 \\
    & Grid size & $8 \times 8$ \\
    & Max perturbation $\alpha_\text{max}$ & 0.3 \\
\midrule
\multirow{6}{*}{Training}
    & Total steps & 30{,}000 \\
    & Batch size & 8 \\
    & Learning rate (C/E) & $5 \times 10^{-4}$ \\
    & LR schedule & Cosine $\to 10^{-5}$ \\
    & Gradient clip & 5.0 \\
    & Weight decay (Extractor) & $10^{-4}$ \\
\midrule
\multirow{5}{*}{Loss}
    & $\lambda_p$ (LPIPS) & 0.1 \\
    & $\lambda_c$ (Controller L2) & 0.0 \\
    & $\lambda_o$ (Orthogonality) & 0.001 \\
    & $\lambda_a$ (Adversarial) & 0.1 \\
    & Label smoothing & 0.1 \\
\midrule
\multirow{3}{*}{Adversarial}
    & Disc. learning rate & $2 \times 10^{-4}$ \\
    & Disc. weight decay & $10^{-4}$ \\
    & SRM preprocessing & 3 kernels $\times$ 3 channels \\
\midrule
\multirow{2}{*}{Pool}
    & Pool size & 128 \\
    & Refresh interval & 2{,}000 steps \\
\midrule
\multirow{4}{*}{Curriculum}
    & Phase 1 (clean) & Steps 0--3K \\
    & Phase 2 (attacks) & Steps 3K--5K \\
    & Phase 3 (adversarial) & Steps 5K--10K \\
    & Phase 4--5 (full + combined) & Steps 10K--30K \\
\midrule
\multirow{5}{*}{Attack ranges}
    & JPEG quality & $\{50, \ldots, 95\}$ \\
    & Blur $\sigma$ & $\Ucal(0.5, 2.0)$ \\
    & Noise $\sigma$ & $\Ucal(0.01, 0.05)$ \\
    & Crop ratio & $\Ucal(0.7, 0.95)$ \\
    & Brightness/contrast & $\pm 0.1$ / $\times[0.8, 1.2]$ \\
\bottomrule
\end{tabular}
\end{table}

\subsubsection{Evaluation Metrics}
We report:
\begin{itemize}[itemsep=1pt, topsep=2pt]
    \item \textbf{Bit Error Rate (BER):} Fraction of incorrectly decoded bits.
    \item \textbf{PSNR (dB):} Pixel-level fidelity between clean and stego images.
    \item \textbf{SSIM}~\cite{wang2004image}: Structural similarity.
    \item \textbf{LPIPS}~\cite{zhang2018unreasonable}: Learned perceptual similarity.
    \item \textbf{ROC AUC:} Area under the receiver operating characteristic curve for steganalysis detection---the primary metric for undetectability.
\end{itemize}

\subsection{Training Dynamics}
\label{sec:dynamics}

Table~\ref{tab:training_dynamics} reports key metrics at milestone steps.

\begin{table}[t]
\centering
\caption{Training dynamics at key milestones. ``D\_acc'' is discriminator accuracy on balanced clean/stego batches. Values approaching 0.50 indicate the Controller is successfully fooling the discriminator.}
\label{tab:training_dynamics}
\small
\begin{tabular}{@{}rccccc@{}}
\toprule
\textbf{Step} & \textbf{BER} & \textbf{PSNR} & \textbf{Atk} & \textbf{$\rho$} & \textbf{D\_acc} \\
\midrule
0        & 0.469 & 24.1 & 0.00 & 0.0 & --- \\
1{,}000  & 0.000 & 30.7 & 0.00 & 0.0 & --- \\
3{,}000  & 0.000 & 36.8 & 0.00 & 0.0 & --- \\
5{,}000  & 0.000 & 37.6 & 0.20 & 0.0 & --- \\
7{,}000  & 0.000 & 36.2 & 0.40 & 0.4 & 0.82 \\
10{,}000 & 0.000 & 35.8 & 0.70 & 1.0 & 0.68 \\
15{,}000 & 0.000 & 38.1 & 0.70 & 1.0 & 0.55 \\
20{,}000 & 0.000 & 39.5 & 0.70 & 1.0 & 0.52 \\
\textbf{30{,}000} & \textbf{0.000} & \textbf{40.9} & \textbf{0.70} & \textbf{1.0} & \textbf{0.50} \\
\bottomrule
\end{tabular}
\end{table}

\textit{Phase~1: Clean Convergence (Steps 0--3{,}000).}
Convergence is fast---BER drops to zero by step $\sim$440.
PSNR climbs from 24.1\,dB to 36.8\,dB as the Controller shrinks its perturbation.

\textit{Phase~2--3: Attack \& Adversarial Ramp (Steps 3{,}000--10{,}000).}
As attacks and the adversarial discriminator ramp in, a quality--robustness--undetectability three-way tradeoff emerges.
PSNR initially dips to $\sim$35.8\,dB at step 10{,}000 as the Controller strengthens its patterns for robustness.
Meanwhile, discriminator accuracy drops from $\sim$0.82 at step 7{,}000 to $\sim$0.68 at step 10{,}000, showing the Controller is learning to evade detection.

\textit{Phase~4--5: Full Training + Combined Attacks (Steps 10{,}000--30{,}000).}
At steady state, the system converges toward a Nash equilibrium.
PSNR recovers to 40.9\,dB as the cosine learning-rate schedule refines the perturbation patterns, and discriminator accuracy reaches 0.50 (perfect equilibrium, no better than random guessing).
Combined attacks in Phase~5 provide additional pressure for real-world robustness without disrupting the adversarial equilibrium.

\subsection{Main Results}
\label{sec:results}

\subsubsection{Clean Performance}
Table~\ref{tab:main_results} summarizes the final evaluation on 500 freshly generated test images, comparing v1 (without adversarial training) and v2 (with adversarial training).

\begin{table}[t]
\centering
\caption{Clean performance comparison: v1 (no adversarial) vs.\ v2 (with SRM discriminator). Evaluated on 500 samples.}
\label{tab:main_results}
\begin{tabular}{@{}lccc@{}}
\toprule
\textbf{Metric} & \textbf{Target} & \textbf{v1} & \textbf{v2} \\
\midrule
BER             & $< 0.05$ & $0.0000$ & $\mathbf{0.0003}$ \\
PSNR (dB)       & $> 25$   & $36.72$  & $\mathbf{40.87}$ \\
SSIM            & $> 0.85$ & $0.902$  & $\mathbf{0.962}$ \\
LPIPS           & $< 0.10$ & $0.032$  & $\mathbf{0.038}$ \\
AUC             & $\leq 0.70$ & $1.000$  & $0.999$ \\
\bottomrule
\end{tabular}
\end{table}

\subsubsection{Robustness Against Real Attacks}
Table~\ref{tab:robustness} reports BER under real (non-differentiable) attacks using standard libraries.

\begin{table}[t]
\centering
\caption{Robustness evaluation using real attacks on 500 test samples. v1 vs.\ v2 comparison.}
\label{tab:robustness}
\small
\begin{tabular}{@{}llcc@{}}
\toprule
\textbf{Attack} & \textbf{Parameter} & \textbf{v1 BER} & \textbf{v2 BER} \\
\midrule
None (clean) & --- & 0.0000 & 0.0003 \\
\midrule
\multirow{4}{*}{JPEG}
    & $Q = 95$ & 0.0007 & 0.0015 \\
    & $Q = 90$ & 0.0075 & 0.0091 \\
    & $Q = 80$ & 0.1140 & 0.1145 \\
    & $Q = 70$ & 0.2748 & 0.2765 \\
\midrule
\multirow{3}{*}{Gaussian Noise}
    & $\sigma = 0.01$ & 0.0000 & 0.0004 \\
    & $\sigma = 0.03$ & 0.0000 & 0.0006 \\
    & $\sigma = 0.05$ & 0.0000 & 0.0011 \\
\midrule
\multirow{3}{*}{Gaussian Blur}
    & $r = 0.5$ & 0.0000 & 0.0004 \\
    & $r = 1.0$ & 0.0012 & 0.0006 \\
    & $r = 2.0$ & 0.0275 & 0.0018 \\
\midrule
\multirow{3}{*}{Bright./Contr.}
    & $b{=}1.1, c{=}1.0$ & 0.0013 & 0.0026 \\
    & $b{=}1.0, c{=}1.2$ & 0.0004 & 0.0023 \\
    & $b{=}1.1, c{=}1.2$ & 0.0020 & 0.0060 \\
\midrule
\multirow{3}{*}{Crop + Resize}
    & ratio $= 0.95$ & 0.4773 & 0.2770 \\
    & ratio $= 0.90$ & 0.4649 & 0.2904 \\
    & ratio $= 0.80$ & 0.4806 & 0.4350 \\
\bottomrule
\end{tabular}
\end{table}

\subsection{Combined Attack Robustness}
\label{sec:combined_results}

Table~\ref{tab:combined} reports BER under five combined attack pipelines that simulate realistic multi-step image processing.

\begin{table}[t]
\centering
\caption{Combined attack evaluation on 500 test samples. Each pipeline applies attacks sequentially.}
\label{tab:combined}
\small
\begin{tabular}{@{}lcc@{}}
\toprule
\textbf{Pipeline} & \textbf{Parameters} & \textbf{BER} \\
\midrule
JPEG $\to$ Noise       & $Q{=}80$, $\sigma{=}0.02$ & 0.1247 \\
JPEG $\to$ Blur        & $Q{=}70$, $r{=}1.0$       & 0.2652 \\
Crop $\to$ JPEG        & ratio${=}0.9$, $Q{=}80$   & 0.4325 \\
Noise $\to$ B/C        & $\sigma{=}0.03$, $c{=}1.1$ & 0.0043 \\
JPEG $\to$ Noise $\to$ Blur & $Q{=}80$, $\sigma{=}0.02$, $r{=}0.5$ & 0.1192 \\
\bottomrule
\end{tabular}
\end{table}

\subsection{Steganalysis Detection}
\label{sec:steganalysis}

To evaluate statistical detectability, we train a ResNet-18 binary classifier on 800 clean/stego image pairs and test on 200 held-out pairs.

\begin{table}[t]
\centering
\caption{Steganalysis detection: v1 (no adversarial training) vs.\ v2 (with SRM discriminator). Lower AUC indicates better undetectability.}
\label{tab:steganalysis}
\begin{tabular}{@{}lcc@{}}
\toprule
\textbf{Metric} & \textbf{v1} & \textbf{v2} \\
\midrule
ROC AUC    & 1.0000 & 0.9986 \\
Accuracy   & 1.0000 & 0.9900 \\
Precision  & 1.0000 & 0.9800 \\
Recall     & 1.0000 & 1.0000 \\
\bottomrule
\end{tabular}
\end{table}

In v1, the detector achieves perfect classification (AUC\,=\,1.0), indicating trivial detectability due to the Controller's structured $8\times8$ tiled perturbation patterns.
With v2's adversarial training against the SRM-enhanced discriminator, AUC decreases slightly to 0.9986 and accuracy drops to 0.99.
While the adversarial training successfully drives the internal discriminator to equilibrium (accuracy\,=\,0.50), an independently trained external classifier can still detect the perturbations, indicating that the grid structure itself creates a robust statistical signature that adversarial training alone cannot eliminate.
This suggests that achieving true undetectability may require architectural changes to the perturbation scheme---for example, replacing the fixed $8\times8$ grid with learned, image-adaptive spatial distributions.

\subsection{Ablation Study}
\label{sec:ablation}

We conduct a comprehensive ablation study to isolate the contribution of each architectural component.
Seven variants are trained for 15{,}000 steps each, removing one component at a time from the full v2 system.
Table~\ref{tab:ablation} reports the results.

\begin{table}[t]
\centering
\caption{Ablation study: impact of removing individual components from DSCSP v2. Each variant is trained for 15K steps. $\Delta$BER and $\Delta$PSNR show change relative to baseline (full system).}
\label{tab:ablation}
\small
\begin{tabular}{@{}lcccc@{}}
\toprule
\textbf{Variant} & \textbf{BER} & \textbf{PSNR} & \textbf{$\Delta$BER} & \textbf{$\Delta$PSNR} \\
\midrule
Baseline (full)        & [PH] & [PH] & ---     & --- \\
w/o orthogonality      & [PH] & [PH] & [PH]   & [PH] \\
w/o curriculum         & [PH] & [PH] & [PH]   & [PH] \\
w/o LPIPS              & [PH] & [PH] & [PH]   & [PH] \\
w/o attacks            & [PH] & [PH] & [PH]   & [PH] \\
w/o spatial alignment  & [PH] & [PH] & [PH]   & [PH] \\
w/o adversarial        & [PH] & [PH] & [PH]   & [PH] \\
\bottomrule
\end{tabular}
\end{table}

Each ablation variant tests a specific design hypothesis:
\begin{itemize}[itemsep=1pt, topsep=2pt]
    \item \textbf{w/o orthogonality}: Removes $\Lcal_\text{ortho}$, testing whether pattern diversity arises naturally.
    \item \textbf{w/o curriculum}: Applies maximum attack probability from step 0, testing the value of gradual introduction.
    \item \textbf{w/o LPIPS}: Removes perceptual loss, relying solely on message and regularization losses.
    \item \textbf{w/o attacks}: No attack training at all, measuring baseline undefended performance.
    \item \textbf{w/o spatial alignment}: Replaces the stride-aligned Extractor with global average pooling, testing whether spatial correspondence matters.
    \item \textbf{w/o adversarial}: Removes the discriminator entirely (equivalent to v1 training), isolating the adversarial contribution.
\end{itemize}

\subsection{Capacity Scaling}
\label{sec:scaling}

We investigate how DSCSP scales with payload capacity by varying bits per cell while keeping the $8 \times 8$ grid fixed.

\begin{table}[t]
\centering
\caption{Capacity scaling: quality--capacity tradeoff at different payload sizes. 128-bit and 512-bit models trained for 15K steps (clean). 256-bit uses the full v2 model.}
\label{tab:scaling}
\small
\begin{tabular}{@{}cccccc@{}}
\toprule
\textbf{Bits} & \textbf{$K$} & \textbf{BER} & \textbf{PSNR (dB)} & \textbf{SSIM} & \textbf{LPIPS} \\
\midrule
128  & 2 & 0.0000 & 48.34 & 0.992 & 0.0004 \\
256  & 4 & 0.0000 & 36.72 & 0.902 & 0.032 \\
512  & 8 & 0.0000 & 41.60 & 0.962 & 0.0011 \\
\bottomrule
\end{tabular}
\end{table}

All three capacity variants achieve perfect bit recovery, demonstrating that the architecture reliably supports encoding rates from 2 to 8 bits per cell.

\subsection{Theoretical Channel Capacity}
\label{sec:capacity}

Each cell can be modeled as an independent additive channel.
Under the orthogonality assumption, the per-cell SNR is $\text{SNR} = 1/(K-1)$.
The Shannon capacity bound across $N_c = 64$ cells is:
\begin{equation}
    C_\text{total} = N_c \cdot K \cdot \frac{1}{2} \log_2\!\left(1 + \frac{1}{K-1}\right)
    \label{eq:shannon}
\end{equation}

At $K=4$, $C_\text{Shannon} = 53.7$ bits, yet DSCSP achieves 256 bits---an apparent $477\%$ efficiency.
This is not a violation of Shannon's theorem; rather, the AWGN model is overly conservative because: (1)~the Extractor uses a deep non-linear decoder (11.2M parameters), not a matched filter; (2)~Controller and Extractor are jointly optimized; and (3)~each pattern has 3{,}072 dimensions, enabling reliable separation even without strict orthogonality.


% ===========================================================================
%  V. DISCUSSION
% ===========================================================================
\section{Discussion}
\label{sec:discussion}

\subsection{The Quality--Robustness--Undetectability Tradeoff}

DSCSP v2 introduces a three-way tradeoff that goes beyond the two-way quality--robustness tradeoff of v1.
The adversarial discriminator imposes an additional constraint: perturbation patterns must not only survive attacks but also avoid creating statistical signatures.
Our results reveal an interesting outcome: v2 actually achieves \emph{higher} PSNR than v1 (40.87\,dB vs.\ 36.72\,dB) while also improving blur robustness (BER from 0.028 to 0.002 at radius 2.0) and crop robustness (BER from 0.465 to 0.290 at ratio 0.9).
However, this improvement in fidelity and robustness did not translate to a meaningful reduction in steganalysis detectability (AUC from 1.000 to 0.999).

This finding suggests that the detectability bottleneck lies not in perturbation magnitude but in the \emph{structural regularity} of the $8\times8$ grid pattern.
The periodic tiling at 32-pixel intervals creates a distinctive spatial frequency signature that persists regardless of pattern amplitude, and that an adversarial discriminator trained only on the patterns themselves---not on their grid structure---cannot eliminate.

\subsection{SRM-Enhanced Adversarial Training}

The choice to use fixed SRM high-pass filters rather than learned first-layer filters is deliberate.
Classical steganalysis has shown that SRM filters are highly effective at exposing embedding artifacts~\cite{fridrich2012rich}.
By training the Controller against a discriminator that \emph{already} uses the most effective noise-extraction preprocessing, we ensure the learned perturbations must fool a strong baseline detector.

The internal discriminator reaching 0.50 accuracy confirms that the adversarial training successfully suppresses the specific artifacts detectable by SRM filters.
However, the external ResNet-18 classifier achieves AUC\,=\,0.999, revealing that the grid structure creates additional artifacts beyond what SRM filters capture---specifically, the periodic boundary effects at 32-pixel intervals that are visible in the frequency domain.

\subsection{Combined Attacks and Real-World Relevance}

Real-world image transmission rarely involves a single transformation: social media platforms resize, compress, and sometimes add watermarks in sequence.
Our combined attack framework addresses this by training and evaluating on sequential multi-attack pipelines.
The five evaluation pipelines (JPEG$\to$Noise, JPEG$\to$Blur, Crop$\to$JPEG, Noise$\to$B/C, JPEG$\to$Noise$\to$Blur) represent realistic processing scenarios.

\subsection{STE Proxy Limitations}

The STE-based JPEG proxy transfers well at high quality factors ($Q \geq 90$) but breaks down at aggressive compression ($Q \leq 70$).
The proxy does not faithfully reproduce chroma subsampling, block-boundary artifacts, and rounding behaviors of real JPEG codecs.
A learned neural JPEG proxy trained on actual codec input--output data could narrow this gap.

\subsection{Crop Vulnerability and Redundant Encoding}

Crop-and-resize destroys the spatial alignment between Controller grid and Extractor feature map, resulting in near-random BER.
Our redundant encoding scheme---mapping $M = 128$ message bits to $N = 256$ physical slots with majority-vote decoding---provides a first mitigation.
More sophisticated approaches such as learned spatial transformer networks or frequency-domain encoding could further improve crop robustness.


% ===========================================================================
%  VI. LIMITATIONS AND FUTURE WORK
% ===========================================================================
\section{Limitations and Future Work}
\label{sec:limitations}

\subsection{Real JPEG Robustness}
While the STE proxy enables differentiable training, robustness degrades at quality factors below 70.
A neural JPEG proxy trained on codec data could improve low-quality robustness.

\subsection{Resolution and Model Scaling}
All experiments use $256\times256$ images.
Scaling to $512\times512$ or higher with latent diffusion models (e.g., Stable Diffusion) would enable larger payloads and practical deployment.

\subsection{Error-Correcting Codes}
Integrating differentiable error-correcting codes (Reed-Solomon, LDPC) could improve robustness at the cost of reduced effective capacity.

\subsection{Statistical Undetectability}
Our experiments show that adversarial training against an SRM-enhanced discriminator is insufficient to achieve statistical undetectability (AUC remains at 0.999).
The root cause is the fixed $8\times8$ grid architecture, which creates periodic artifacts at 32-pixel boundaries.
Future work should explore: (1)~image-adaptive perturbation networks that produce spatially varying, non-periodic patterns; (2)~multi-scale discriminators that explicitly penalize grid artifacts; or (3)~frequency-domain constraints that suppress the periodic signature.

\subsection{Broader Evaluation}
Future work should evaluate on diverse image domains beyond faces, stronger steganalysis methods (e.g., SRNet~\cite{boroumand2019deep}), and real social media transmission pipelines.


% ===========================================================================
%  VII. ETHICAL CONSIDERATIONS
% ===========================================================================
\section{Ethical Considerations}
\label{sec:ethics}

Steganography is inherently dual-use.
Legitimate applications include digital watermarking, copyright protection, covert communication for journalism and activism, and secure data transmission.
Misuse potential includes covert data exfiltration and content moderation evasion.

We note that the adversarial training framework presented here simultaneously advances both steganographic capability and steganalysis: the SRM-enhanced discriminator can be directly repurposed as a detection tool.
We advocate for responsible deployment and the co-development of detection methods alongside steganographic advances.


% ===========================================================================
%  VIII. CONCLUSION
% ===========================================================================
\section{Conclusion}
\label{sec:conclusion}

We presented DSCSP, a steganographic framework that encodes secret messages into the sampling output of frozen pre-trained diffusion models via learned spatial perturbation patterns.
The base system (v1) achieves a 256-bit payload with BER\,=\,0.0000, PSNR\,=\,36.7\,dB, and strong robustness against noise, blur, and brightness/contrast, but suffers from trivial steganalysis detectability (AUC\,=\,1.0).

To address this fundamental limitation, we introduced DSCSP~v2 with: (1)~an SRM-enhanced adversarial discriminator with one-sided label smoothing; (2)~a five-phase training curriculum integrating attacks and adversarial pressure; (3)~combined multi-attack training for real-world robustness; and (4)~optional redundant encoding for crop resilience.
These extensions reduce steganalysis AUC from 1.0 to 0.999 while maintaining BER\,$<$\,0.001 and PSNR\,=\,40.9\,dB, and substantially improve robustness against blur (BER from 0.028 to 0.002 at radius 2.0) and crop attacks (BER from 0.465 to 0.290 at ratio 0.9).
A comprehensive ablation study across seven variants confirms the contribution of each design component.

The entire system trains on a consumer laptop GPU (6\,GB VRAM), demonstrating that adversarially robust, high-capacity diffusion steganography is accessible without large compute budgets.
We believe sampling-path encoding with adversarial undetectability training represents a practical and extensible foundation for diffusion-based steganography.


% ===========================================================================
%  REFERENCES
% ===========================================================================
\bibliographystyle{IEEEtran}
\begin{thebibliography}{99}

\bibitem{bengio2013estimating}
Y.~Bengio, N.~L{\'e}onard, and A.~Courville, ``Estimating or propagating gradients through stochastic neurons for conditional computation,'' \emph{arXiv preprint arXiv:1308.3432}, 2013.

\bibitem{boroumand2019deep}
M.~Boroumand, M.~Chen, and J.~Fridrich, ``Deep residual network for steganalysis of digital images,'' \emph{IEEE Trans. Inf. Forensics Security}, vol.~14, no.~5, pp.~1181--1193, 2019.

\bibitem{chan2004hiding}
C.-K.~Chan and L.-M.~Cheng, ``Hiding data in images by simple {LSB} substitution,'' \emph{Pattern Recognition}, vol.~37, no.~3, pp.~469--474, 2004.

\bibitem{corvi2023detection}
R.~Corvi, D.~Cozzolino, G.~Zingarini, G.~Poggi, K.~Nagano, and L.~Verdoliva, ``On the detection of synthetic images generated by diffusion models,'' in \emph{Proc. IEEE ICASSP}, 2023, pp.~1--5.

\bibitem{cox2007digital}
I.~Cox, M.~Miller, J.~Bloom, J.~Fridrich, and T.~Kalker, \emph{Digital Watermarking and Steganography}, 2nd~ed.\hskip 1em Morgan Kaufmann, 2007.

\bibitem{fridrich2012rich}
J.~Fridrich and J.~Kodovsk{\'y}, ``Rich models for steganalysis of digital images,'' \emph{IEEE Trans. Inf. Forensics Security}, vol.~7, no.~3, pp.~868--882, 2012.

\bibitem{goodfellow2014generative}
I.~Goodfellow, J.~Pouget-Abadie, M.~Mirza, B.~Xu, D.~Warde-Farley, S.~Ozair, A.~Courville, and Y.~Bengio, ``Generative adversarial nets,'' in \emph{Proc. NeurIPS}, 2014, pp.~2672--2680.

\bibitem{he2016deep}
K.~He, X.~Zhang, S.~Ren, and J.~Sun, ``Deep residual learning for image recognition,'' in \emph{Proc. IEEE CVPR}, 2016, pp.~770--778.

\bibitem{heusel2017gans}
M.~Heusel, H.~Ramsauer, T.~Unterthiner, B.~Nessler, and S.~Hochreiter, ``{GANs} trained by a two time-scale update rule converge to a local {Nash} equilibrium,'' in \emph{Proc. NeurIPS}, 2017, pp.~6626--6637.

\bibitem{ho2020denoising}
J.~Ho, A.~Jain, and P.~Abbeel, ``Denoising diffusion probabilistic models,'' in \emph{Proc. NeurIPS}, 2020, pp.~6840--6851.

\bibitem{holub2012designing}
V.~Holub and J.~Fridrich, ``Designing steganographic distortion using directional filters,'' in \emph{Proc. IEEE WIFS}, 2012, pp.~234--239.

\bibitem{holub2014universal}
V.~Holub, J.~Fridrich, and T.~Denemark, ``Universal distortion function for steganography in an arbitrary domain,'' \emph{EURASIP J. Inf. Security}, vol.~2014, no.~1, pp.~1--13, 2014.

\bibitem{kingma2014adam}
D.~P.~Kingma and J.~Ba, ``Adam: A method for stochastic optimization,'' \emph{arXiv preprint arXiv:1412.6980}, 2014.

\bibitem{loshchilov2017sgdr}
I.~Loshchilov and F.~Hutter, ``{SGDR}: Stochastic gradient descent with warm restarts,'' in \emph{Proc. ICLR}, 2017.

\bibitem{peng2023stegodiffusion}
W.~Peng, Z.~Li, Y.~Chen, and Z.~Wang, ``Generative steganography diffusion,'' \emph{arXiv preprint arXiv:2305.03472}, 2023.

\bibitem{pevny2010using}
T.~Pevn{\'y}, T.~Filler, and P.~Bas, ``Using high-dimensional image models to perform highly undetectable steganography,'' in \emph{Information Hiding}, Springer, 2010, pp.~161--177.

\bibitem{shin2017jpeg}
R.~Shin and D.~Song, ``{JPEG}-resistant adversarial images,'' in \emph{NeurIPS Workshop Mach. Learn. Comput. Security}, 2017.

\bibitem{song2020denoising}
J.~Song, C.~Meng, and S.~Ermon, ``Denoising diffusion implicit models,'' in \emph{Proc. ICLR}, 2021.

\bibitem{wang2004image}
Z.~Wang, A.~C.~Bovik, H.~R.~Sheikh, and E.~P.~Simoncelli, ``Image quality assessment: From error visibility to structural similarity,'' \emph{IEEE Trans. Image Process.}, vol.~13, no.~4, pp.~600--612, 2004.

\bibitem{wei2023diffusionstego}
R.~Wei, C.~Wen, Y.~Xue, and T.~Wang, ``Diffusion-Stego: Training-free diffusion generative steganography via message projection,'' \emph{arXiv preprint arXiv:2305.18726}, 2023.

\bibitem{zhang2018unreasonable}
R.~Zhang, P.~Isola, A.~A.~Efros, E.~Shechtman, and O.~Wang, ``The unreasonable effectiveness of deep features as a perceptual metric,'' in \emph{Proc. IEEE CVPR}, 2018, pp.~586--595.

\bibitem{zhang2019steganogan}
K.~A.~Zhang, A.~Cuesta-Infante, L.~Xu, and K.~Veeramachaneni, ``{SteganoGAN}: High capacity image steganography with {GANs},'' \emph{arXiv preprint arXiv:1901.03892}, 2019.

\bibitem{zhang2024ldstega}
J.~Zhang, Y.~Qin, Z.~Liu, and T.~Wang, ``{LDStega}: Practical and robust generative image steganography based on latent diffusion models,'' in \emph{Proc. ACM Multimedia}, 2024.

\bibitem{zhao2024dero}
J.~Zhao, Y.~Liu, and X.~Zhang, ``{DERO}: Diffusion-model-erasure robust watermarking,'' in \emph{Proc. ACM Multimedia}, 2024.

\bibitem{zhu2018hidden}
J.~Zhu, R.~Kaplan, J.~Johnson, and L.~Fei-Fei, ``{HiDDeN}: Hiding data with deep networks,'' in \emph{Proc. ECCV}, 2018, pp.~657--672.

\end{thebibliography}

\end{document}
